\label {ch:model_and_methods}

\section{Постановка ресурсной СМО}
\label{sec:formal_model}

В данном разделе будет смоделирована и рассмотрена ресурсная loss система массового обслуживания, где каждая поступающая заявка занимает 1 канал и случайный объем ограниченного ресурса, а скорость поступления заявок и время их обработки зависит постоянны только от характера их распределения. 

Основные параметры модели:
\begin{table}[h!]
    \centering
    \caption{Основные параметры модели}
    \begin{tabular}{|p{7cm}|l|}
        \hline
        \textbf{Наименование параметра} & \textbf{Обозначение} \\ \hline
        Количество приборов & $N = \text{SERVERS\_N}$ \\ \hline
        Число обслуживающих каналов & $K = N$ \\ \hline
        Вместимость системы по ресурсам & $Res = \text{TOTAL\_RESOURCE\_R}$ \\ \hline
        Средняя интенсивность поступления & $\lambda$ \\ \hline
        Скорость обслуживания & $\sigma$ \\ \hline
        Случайное требование заявки к ресурсу & $D$ \\ \hline
        Случайный объем работы заявки & $B$ \\ \hline
        Фактическое время работы заявки & $S = \frac{B}{\sigma}$ \\ \hline
    \end{tabular}
\end{table}

Состояние системы в агрегированном виде описываются через $n(t)$ и $u(t)$ - число заявок и занятый ресурс в момент времени $t$. Новая заявка при попадании в систему имеет $D \sim F_D$ и $B \sim F_B$. Заявка принимается, если одновременно выполнены ограничения:
%\usepackage{amsmath} 
$
\begin{cases}
    n(t) < N \\
    u(t) + D \leq R
\end{cases}
$

Так как система не имеет буфера, то заявка теряется, если не выполнено хотя бы одно ограничение. Заявка обслуживается $S=\frac{B}{\sigma}$, после чего освобождает ресурс и канал. 

В работе рассматривается ресурсную многоканальную loss-СМО с одним типом ограниченного ресурса, случайными ресурсными требованиями, общим распределением объема работы и варьируемым типом входящего потока $GI/G/N/N/R$. $GI$ отвечает за тип входящего потока, $G$ --- за распределение объема работы заявки, $N$ --- за число обслуживающих приборов и за число заявок в системе, $R$ --- за общий объем ресурса в системе.

Такая постановка задачи является общим случаем $M/G/N/N/R$ системы. Система не являетс марковской, потому что на вход подается не только экспоненциально распределенный поток заявок, ведь для предсказания будущего состояния необходимо знать остаточное время до следующего поступления. Формально систему можно сделать марковской, если расширить состояние, храня фазы поступлений и обслуживаний, однако это приведет к резкому росту размерности и компьютерная имитация становится слишком затратной. 

Таким образом агрегированная модель $(n(t),u(t))$ не является марковской при общем входящем потоке и неэкспоненциальном workload. Отдельные частные случаи сочетаний распределений входящего потока и времени работы будут рассмотрены в 3й главе. 

\section{Исследуемые метрики и показатели чувствительности}
\label{sec:metrics}

Каждый сценарий имитационного эксперимента задаётся набором параметров
$$
s = (a,\omega,\lambda,\sigma,N,R,F_D),
$$
где $a$ --- тип входящего потока, $\omega$ --- распределение объёма работы заявки, $\lambda$ --- средняя интенсивность поступления, $\sigma$ --- скорость обслуживания, $N$ --- число обслуживающих каналов, $R$ --- общий объём ресурса, $F_D$ --- распределение ресурсного требования.

Результатом моделирования сценария $s$ является набор функционалов системы $m = m(s)$.
В работе анализируются следующие основные метрики: вероятность отказа, пропускная способность, среднее число заявок в системе, средний занятый ресурс, среднее время обслуживания и среднее время пребывания заявки в системе.

Вероятность отказа оценивается как доля потерянных заявок среди всех поступивших:
$$
\widehat{P}_{\mathrm{loss}}
=
\frac{N_{\mathrm{rej}}}{N_{\mathrm{arr}}},
$$
где $N_{\mathrm{arr}}$ --- число поступивших заявок, $N_{\mathrm{rej}}$ --- число заявок, отклонённых из-за нарушения ограничений по числу серверов или доступному ресурсу.

Пропускная способность системы определяется как среднее число успешно принятых и завершённых заявок в единицу времени наблюдения:
$$
\widehat{\theta}
=
\frac{N_{\mathrm{acc}}}{T_{\mathrm{obs}}},
$$
где $N_{\mathrm{acc}}$ --- число принятых заявок, $T_{\mathrm{obs}} = T_{\max} - T_{\mathrm{warm}}$ --- длина интервала, на котором собирается статистика. 

Среднее число заявок в системе оценивается как временное среднее:
$$
\widehat{\mathbb{E}}[n]=\frac{1}{T_{\mathrm{obs}}} int_{T_{\mathrm{warm}}}^{T_{\max}} n(t)\,dt
$$
Средний занятый ресурс:
$$
\widehat{\mathbb{E}}[u]=\frac{1}{T_{\mathrm{obs}}}\int_{T_{\mathrm{warm}}}^{T_{\max}} u(t)\,dt
$$

Время пребывания $T_{\mathrm{soj}}$ принятой заявки совпадает с её временем обслуживания $S$ и равно $\frac{W}{\sigma}$.

Для анализа чувствительности используется сравнение метрик между базовым и изменённым сценарием. Вероятности отказа оценивает абсолютное отклонение в процентных пунктах:
$$
\Delta P_{\mathrm{loss}}(s,s_0)=100\cdot\left(\widehat{P}_{\mathrm{loss}}(s)-\widehat{P}_{\mathrm{loss}}(s_0)\right).
$$
Для остальных метрик используется относительное отклонение:
$$
\delta_m(s,s_0)=100\cdot\left(\frac{\widehat{m}(s)}{\widehat{m}(s_0)}-1\right).
$$

Базовыми сценариями, относительно которых проверяется смещение результата является детерминированное обслуживание workload и пуассоновский входящий поток заявок.

Пусть для фиксированного сценария $s$ и метрики $m$ выполнено $R_{\mathrm{rep}}$ независимых репликаций. Тогда формируется выборка
$$
Y_{s,m}^{(1)},Y_{s,m}^{(2)},\dots,Y_{s,m}^{(R_{\mathrm{rep}})}.
$$
Точечная оценка метрики вычисляется как выборочное среднее:
$$
\overline{Y}_{s,m}
=
\frac{1}{R_{\mathrm{rep}}}
\sum_{r=1}^{R_{\mathrm{rep}}}Y_{s,m}^{(r)}.
$$
Выборочная дисперсия равна
$$
S_{s,m}^2
=
\frac{1}{R_{\mathrm{rep}}-1}
\sum_{r=1}^{R_{\mathrm{rep}}}
\left(
Y_{s,m}^{(r)}-\overline{Y}_{s,m}
\right)^2.
$$
Стандартная ошибка среднего:
$$
\mathrm{SE}_{s,m}
=
\frac{S_{s,m}}{\sqrt{R_{\mathrm{rep}}}}.
$$
Для оценки статистической надёжности результатов используются доверительные интервалы вида: 
$$
\overline{Y}_{s,m}
\pm
t_{1-\alpha/2,\,R_{\mathrm{rep}}-1}
\cdot
\mathrm{SE}_{s,m},
$$
где $t_{1-\alpha/2,\,R_{\mathrm{rep}}-1}$ --- квантиль распределения Стьюдента. В работе используется уровень доверия $95\%$. 

В рамках работы оценивается абсолютный и относитлеьный сдвиг следующих метрик: вероятности отказа, пропускной способности, среднее число заявок в системе, средний занятый ресурс, среднее время обслуживания и пребывания, декомпозиция отказов. 

\section{Алгоритм имитационного моделирования}
\label{sec:simulation_algorithm}

Для выбранной постановки задачи и удобства компьютерной имитации выбрано моделирование методом Монте-Карло, так как в $GI/G/N/N/R$ системе каждая репликация является независимой реализацией одного и того же стохастического эксперимента.

Статистические характеристики модели оцениваются за время прогона $T_{\max}$, из которого исключается время разогрева $T_{\text{warm}}$. Таким образом, интервал наблюдения имеет длину
$$
T_{\text{obs}} = T_{\max} - T_{\text{warm}}.
$$
Время разогрева необходимо для уменьшения влияния начального состояния системы на стационарные оценки.

В рамках метода Монте-Карло выполняется серия из $R_{\text{rep}}$ независимых репликаций, каждая из которых является одной полной реализацией эксперимента. При фиксированном сценарии $s$ и выбранной целевой метрике $m$ формируется выборка независимых значений
$$
Y_{s,m}^{(1)},Y_{s,m}^{(2)},\dots,Y_{s,m}^{(R_{\text{rep}})}.
$$

Один сценарий задаётся набором параметров
$$
s = (a,\omega,\lambda,\sigma,N,K,R,F_D),
$$
где $a$ --- тип входящего потока, $\omega$ --- распределение объёма работы заявки, $\lambda$ --- средняя интенсивность поступления, $\sigma$ --- скорость обслуживания, $N$ --- число серверов, $K$ --- вместимость системы по числу заявок, $R$ --- общий объём ресурса, $F_D$ --- распределение ресурсного требования. В рассматриваемой loss-системе без очереди выполняется $K=N$.

В начале каждой репликации система находится в пустом состоянии:
$$
n(0)=0,
\qquad
u(0)=0,
$$
где $n(t)$ --- число заявок в системе, $u(t)$ --- занятый ресурс в момент времени $t$.

Входящий поток задаётся последовательностью межприходных интервалов
$$
A_1,A_2,\dots
$$
с распределением, соответствующим выбранному типу arrival-процесса. Моменты поступления заявок определяются рекуррентно:
$$
t_i = t_{i-1} + A_i.
$$

Для каждой поступающей заявки генерируются ресурсное требование
$$
D_i \sim F_D
$$
и объём работы
$$
W_i \sim F_W.
$$
Фактическое время обслуживания определяется формулой
$$
S_i = \frac{W_i}{\sigma}.
$$

Имитация является дискретно-событийной. В каждый момент времени выбирается ближайшее событие: поступление новой заявки или завершение обслуживания одной из уже принятых заявок. Если множество активных заявок обозначить через $\mathcal{I}(t)$, то занятый ресурс имеет вид
$$
u(t)=\sum_{i\in\mathcal{I}(t)}D_i.
$$

При поступлении заявки в момент $t_i$ проверяются ограничения системы:
$$
n(t_i-) < K,
\qquad
u(t_i-) + D_i \leq R.
$$
Если оба условия выполнены, заявка принимается:
$$
n(t_i+) = n(t_i-) + 1,
\qquad
u(t_i+) = u(t_i-) + D_i.
$$
Для принятой заявки планируется момент завершения обслуживания:
$$
c_i = t_i + S_i.
$$

Если хотя бы одно из ограничений нарушено, заявка теряется. В этом случае состояние системы не меняется, но увеличивается счётчик отказов. Отказы дополнительно классифицируются по причине: превышение вместимости $K$, отсутствие доступного сервера $N$ или недостаток ресурса $R$. В текущей loss-постановке $K=N$, однако раздельный учёт причин отказа сохраняется для диагностики.

При завершении обслуживания в момент $c_i$ заявка покидает систему и освобождает один сервер и занятый ресурс:
$$
n(c_i+) = n(c_i-) - 1,
\qquad
u(c_i+) = u(c_i-) - D_i.
$$

Статистика собирается только на интервале
$$
[T_{\mathrm{warm}},T_{\max}].
$$
Если событие происходит до $T_{\mathrm{warm}}$, оно влияет на состояние системы, но не входит в накопление статистики. Это позволяет использовать начальный участок траектории только для выхода системы из пустого состояния.

Между двумя соседними событиями значения $n(t)$ и $u(t)$ постоянны. Поэтому временные средние вычисляются как суммы площадей по межсобытийным интервалам:
$$
\widehat{\mathbb{E}}[n]
=
\frac{1}{T_{\mathrm{obs}}}
\sum_j n_j \Delta t_j,
$$
$$
\widehat{\mathbb{E}}[u]
=
\frac{1}{T_{\mathrm{obs}}}
\sum_j u_j \Delta t_j,
$$
где $\Delta t_j$ --- длина очередного интервала между событиями после учёта разогрева, $n_j$ и $u_j$ --- значения состояния на этом интервале.

Для оценки стационарного распределения числа заявок в системе накапливается суммарное время пребывания в каждом состоянии:
$$
\widehat{\pi}(k)
=
\frac{T_k}{T_{\mathrm{obs}}},
$$
где $T_k$ --- суммарное время, в течение которого после периода разогрева в системе находилось $k$ заявок. Тогда среднее число заявок может быть вычислено как
$$
\widehat{\mathbb{E}}[n]
=
\sum_{k=0}^{K} k\widehat{\pi}(k).
$$

Вероятность отказа оценивается как доля потерянных заявок:
$$
\widehat{P}_{\mathrm{loss}}
=
\frac{N_{\mathrm{rej}}}{N_{\mathrm{arr}}},
$$
где $N_{\mathrm{arr}}$ --- число попыток поступления, $N_{\mathrm{rej}}$ --- число отказов на интервале наблюдения.

Пропускная способность системы оценивается как число завершённых обслуживаний в единицу времени:
$$
\widehat{\theta}
=
\frac{N_{\mathrm{completed}}}{T_{\mathrm{obs}}}.
$$
Для достаточно длинного интервала наблюдения в loss-системе также выполняется приближённая связь
$$
\widehat{\theta}
\approx
\lambda(1-\widehat{P}_{\mathrm{loss}}),
$$
так как доля принятых заявок равна $1-\widehat{P}_{\mathrm{loss}}$.

Средний занятый ресурс вычисляется через накопленное ресурсное время:
$$
\widehat{\mathbb{E}}[u]
=
\frac{\mathrm{resource\_time}}{T_{\mathrm{obs}}}.
$$

Среднее время обслуживания и среднее время пребывания считаются по завершённым заявкам:
$$
\widehat{\mathbb{E}}[S]
=
\frac{1}{N_{\mathrm{completed}}}
\sum_{i=1}^{N_{\mathrm{completed}}} S_i,
$$
$$
\widehat{\mathbb{E}}[T_{\mathrm{soj}}]
=
\frac{1}{N_{\mathrm{completed}}}
\sum_{i=1}^{N_{\mathrm{completed}}} T_{\mathrm{soj},i}.
$$
В рассматриваемой loss-системе без очереди заявка не ожидает начала обслуживания, поэтому для принятой заявки
$$
T_{\mathrm{soj},i}=S_i.
$$
Обе величины сохраняются отдельно, так как в более общей системе с очередью они уже не совпадают.

После выполнения всех репликаций для каждого сценария $s$ получается выборка значений каждой метрики:
$$
Y_{s,m}^{(1)},Y_{s,m}^{(2)},\dots,Y_{s,m}^{(R_{\mathrm{rep}})}.
$$
Точечная оценка метрики равна выборочному среднему:
$$
\overline{Y}_{s,m}
=
\frac{1}{R_{\mathrm{rep}}}
\sum_{r=1}^{R_{\mathrm{rep}}}Y_{s,m}^{(r)}.
$$

Выборочное стандартное отклонение определяется как
$$
S_{s,m}
=
\sqrt{
\frac{1}{R_{\mathrm{rep}}-1}
\sum_{r=1}^{R_{\mathrm{rep}}}
\left(
Y_{s,m}^{(r)}-\overline{Y}_{s,m}
\right)^2
}.
$$

Стандартная ошибка среднего:
$$
\mathrm{SE}_{s,m}
=
\frac{S_{s,m}}{\sqrt{R_{\mathrm{rep}}}}.
$$

Для оценки статистической неопределённости строится доверительный интервал:
$$
\overline{Y}_{s,m}
\pm
z_{1-\alpha/2}\mathrm{SE}_{s,m},
$$
где $z_{1-\alpha/2}$ --- квантиль стандартного нормального распределения. В работе используется уровень доверия $95\%$.

Шум Монте-Карло проявляется в разбросе значений $Y_{s,m}^{(r)}$ между независимыми репликациями. Если различие между двумя сценариями меньше или сопоставимо со стандартной ошибкой, оно не интерпретируется как устойчивый эффект чувствительности. Если же смещение между сценариями существенно превышает внутрисценарный разброс, то оно рассматривается как значимое влияние изменяемого фактора.

Для сравнения двух сценариев $s_1$ и $s_0$ по метрике $m$ используется абсолютное смещение
$$
\Delta_m(s_1,s_0)
=
\overline{Y}_{s_1,m}
-
\overline{Y}_{s_0,m}.
$$
Для вероятности отказа такое смещение удобно выражать в процентных пунктах:
$$
\Delta P_{\mathrm{loss}}(s_1,s_0)
=
100\cdot
\left(
\overline{Y}_{s_1,P_{\mathrm{loss}}}
-
\overline{Y}_{s_0,P_{\mathrm{loss}}}
\right).
$$
Для положительных метрик, таких как пропускная способность или средний занятый ресурс, дополнительно используется относительное смещение:
$$
\delta_m(s_1,s_0)
=
100\cdot
\left(
\frac{\overline{Y}_{s_1,m}}
{\overline{Y}_{s_0,m}}
-
1
\right).
$$

Таким образом, имитационный алгоритм состоит из трёх уровней. На нижнем уровне выполняется дискретно-событийная симуляция одной репликации. На среднем уровне для фиксированного сценария усредняются независимые репликации. На верхнем уровне сравниваются разные сценарии, отличающиеся типом входящего потока, распределением workload и уровнем нагрузки.

Для выбранной постановки задачи и удобства компьютерной имитации выбранно моделирование Монте-Карло, так как в $GI/G/N/N/R$ системе каждая репликация является независимой реализацией одного и того же стохастического эксперимента.

Статистические характеристики модели оцениваются за время прогона $T_{\max}$ из которого вычитается время разогрева $T_{\text{warm}}$; таким образом получаем время наблюдения $T_{\text{obs}}=T_{\max} - T_{\text{warm}}$. Время разогрева необходимо, чтобы нивелировать эффект начального состояния на статистику. 

В рамках метода Монте-Карло выполняется серия из $R\_{\text{rep}}$ независимых репликаций (прогонов), каждая из которых является одной полной реализацией эксперимента. При фиксированном сценарии $s$ и выбранной целевой метрике $m$ формируется выборка независимых значений $Y_{s,m}^{(1)}, Y_{s,m}^{(2)}, \dots, Y_{s,m}^{(R_{\text{rep}})}$. 

На основе полученной выборки далее вычисляются точечные интервальные статистические оценки

События бывают двух типов: поступление новой заявки и завершение обслуживания принятой заявки. 

При поступлении заявки генерируются ресурсное требование $D_i$ и workload $W_i$. Время обслуживания определяется как
$$
S_i = \frac{W_i}{\sigma}.
$$
Заявка принимается, если одновременно выполнены ограничения
$$
n(t_i-) < N,
\qquad
u(t_i-) + D_i \leq R.
$$
Если хотя бы одно ограничение нарушено, заявка теряется. При принятии заявки увеличиваются счётчики принятых заявок, число занятых серверов и занятый ресурс:
$$
n(t_i+) = n(t_i-) + 1,
\qquad
u(t_i+) = u(t_i-) + D_i.
$$
Также планируется момент завершения обслуживания:
$$
c_i = t_i + S_i.
$$

При завершении обслуживания заявка покидает систему, освобождая один сервер и занятый ресурс:
$$
n(c_i+) = n(c_i-) - 1,
\qquad
u(c_i+) = u(c_i-) - D_i.
$$

Статистика собирается только на интервале после разогрева:
$$
[T_{\mathrm{warm}},T_{\max}].
$$
Так как между соседними событиями состояние системы постоянно, временные средние вычисляются через накопление площадей:
$$
\widehat{\mathbb{E}}[n]
=
\frac{1}{T_{\mathrm{obs}}}
\sum_j n_j \Delta t_j,
$$
$$
\widehat{\mathbb{E}}[u]
=
\frac{1}{T_{\mathrm{obs}}}
\sum_j u_j \Delta t_j.
$$
Аналогично для оценки стационарного распределения числа заявок накапливается время пребывания системы в каждом состоянии:
$$
\hat{\pi}(k)
=
\frac{T_k}{T_{\mathrm{obs}}},
$$
где $T_k$ --- суммарное время, в течение которого после периода разогрева в системе находилось $k$ заявок.

$$
\widehat{P}_{\mathrm{loss}}
=
\frac{N_{\mathrm{rej}}}{N_{\mathrm{arr}}},
$$
$$
\widehat{\theta}
=
\frac{N_{\mathrm{completed}}}{T_{\mathrm{obs}}},
$$
$$
\widehat{\mathbb{E}}[u]
=
\frac{\text{resource\_time}}{T_{\mathrm{obs}}},
$$
$$
\widehat{\mathbb{E}}[n]
=
\sum_{k=0}^{K} k\hat{\pi}(k).
$$
Среднее время обслуживания и среднее время пребывания считаются по завершённым заявкам:
$$
\widehat{\mathbb{E}}[S]
=
\frac{1}{N_{\mathrm{completed}}}
\sum_{i=1}^{N_{\mathrm{completed}}} S_i,
$$
$$
\widehat{\mathbb{E}}[T_{\mathrm{soj}}]
=
\frac{1}{N_{\mathrm{completed}}}
\sum_{i=1}^{N_{\mathrm{completed}}} T_{\mathrm{soj},i}.
$$
В рассматриваемой loss-системе без очереди выполняется $T_{\mathrm{soj},i}=S_i$, однако обе величины сохраняются отдельно для унификации формата результатов.

\section{Семейства распределений обслуживания и входного потока}
\label{sec:distributions}

В ресурсной loss-СМО варьируются распределение межприходных интервалов входящего потока и распределение объёма работы заявки, все распределения внутри одного семейства нормируются по математическому ожиданию. 

Пусть $A$ --- случайный интервал между соседними поступлениями. Для всех входящих потоков выполняется условие
$$    
\mathbb{E}[A]=\frac{1}{\lambda}.
$$
Вариативность межприходных интервалов характеризуется квадратом коэффициента вариации
$$
c_A^2=\frac{\operatorname{Var}(A)}{\mathbb{E}[A]^2}.
$$

Пуассоновский входящий поток задаётся экспоненциальным распределением межприходных интервалов:
$$
A \sim \operatorname{Exp}(\lambda),
\qquad
f_A(t)=\lambda e^{-\lambda t}, \quad t \geq 0.
$$
Для него
$$
c_A^2=1.
$$
Этот поток используется как базовый случай, поскольку он обладает марковским свойством и соответствует классической постановке $M/G/N/N/R$.

Регулярные входящие потоки моделируются распределением Эрланга порядка $k$:
$$
A \sim \operatorname{Erlang}(k,k\lambda),
$$
$$
f_A(t)=\frac{(k\lambda)^k t^{k-1}e^{-k\lambda t}}{(k-1)!}, \quad t \geq 0.
$$
При такой параметризации сохраняется среднее значение:
$$
\mathbb{E}[A]=\frac{1}{\lambda},
$$
а вариативность равна
$$
c_A^2=\frac{1}{k}.
$$
В работе используются $k=2$ и $k=4$. Эти потоки задают более равномерное поступление заявок по сравнению с пуассоновским потоком.

Нерегулярный входящий поток моделируется гиперэкспоненциальным распределением второго порядка:
$$
A \sim 
\begin{cases}
\operatorname{Exp}(\mu_1), & \text{с вероятностью } p,\\
\operatorname{Exp}(\mu_2), & \text{с вероятностью } 1-p.
\end{cases}
$$
Его плотность имеет вид
$$
f_A(t)=p\mu_1 e^{-\mu_1 t}+(1-p)\mu_2 e^{-\mu_2 t}.
$$
Параметры $\mu_1$, $\mu_2$ и $p$ выбираются так, чтобы выполнялось условие
$$
\mathbb{E}[A]=\frac{p}{\mu_1}+\frac{1-p}{\mu_2}=\frac{1}{\lambda}.
$$
Для гиперэкспоненциального распределения характерно
$$
c_A^2>1.
$$
Поэтому такой входящий поток используется для моделирования bursty-нагрузки: интервалы между заявками чередуют короткие плотные участки и более длинные паузы.

Второй варьируемый механизм --- распределение объёма работы заявки. Пусть $W$ --- случайный объём работы, а фактическое время обслуживания равно
$$
S=\frac{W}{\sigma}.
$$
Во всех сценариях распределения workload нормируются так, что
$$
\mathbb{E}[W]=1,
\qquad
\mathbb{E}[S]=\frac{1}{\sigma}.
$$
Следовательно, сравнение workload-семейств проводится при одинаковом среднем времени обслуживания, но при разной вариативности:
$$
c_W^2=\frac{\operatorname{Var}(W)}{\mathbb{E}[W]^2}.
$$

Детерминированный workload задаётся как
$$
W \equiv 1,
\qquad
c_W^2=0.
$$
Он соответствует полностью предсказуемому времени обслуживания.

Экспоненциальный workload задаётся распределением
$$
W \sim \operatorname{Exp}(1),
\qquad
c_W^2=1.
$$
Этот случай используется как классический марковский базис по обслуживанию.

Слабовариативный workload моделируется распределением Эрланга:
$$
W \sim \operatorname{Erlang}(k,k),
$$
где в экспериментах используется $k=4$. Тогда
$$
\mathbb{E}[W]=1,
\qquad
c_W^2=\frac{1}{4}.
$$

Высоковариативный workload задаётся гиперэкспоненциальным распределением:
$$
W \sim
\begin{cases}
\operatorname{Exp}(\nu_1), & \text{с вероятностью } q,\\
\operatorname{Exp}(\nu_2), & \text{с вероятностью } 1-q.
\end{cases}
$$
Параметры выбираются из условия
$$
\frac{q}{\nu_1}+\frac{1-q}{\nu_2}=1.
$$
Такое распределение имеет $c_W^2>1$ и моделирует наличие коротких задач вместе с редкими длительными заявками.

Таким образом, в работе сравниваются распределения с одинаковыми первыми моментами, но с различными коэффициентами вариации. Это позволяет отдельно исследовать влияние формы входящего потока и распределения workload на характеристики ресурсной системы.

% \section{Исследуемые метрики и показатели чувствительности}
% \label{sec:metrics}

% Каждый сценарий имитационного эксперимента задаётся набором параметров
% $$
% s=(a,\omega,\lambda,\sigma,N,K,R,F_D),
% $$
% где $a$ --- тип входящего потока, $\omega$ --- распределение workload, $\lambda$ --- средняя интенсивность поступления, $\sigma$ --- скорость обслуживания, $N$ --- число серверов, $K$ --- вместимость системы по числу заявок, $R$ --- общий объём ресурса, $F_D$ --- распределение ресурсного требования. В данной loss-постановке
% $$
% K=N.
% $$

% Результатом моделирования сценария $s$ является набор функционалов системы
% $$
% m=m(s).
% $$
% В работе используются следующие метрики:
% $$
% P_{\mathrm{loss}},\quad
% \theta,\quad
% \mathbb{E}[n],\quad
% \mathbb{E}[u],\quad
% \mathbb{E}[S],\quad
% \mathbb{E}[T_{\mathrm{soj}}].
% $$
% Они описывают соответственно вероятность отказа, пропускную способность, среднее число заявок в системе, средний занятый ресурс, среднее время обслуживания и среднее время пребывания заявки в системе.

% Вероятность отказа оценивается как доля потерянных заявок:
% $$
% \widehat{P}_{\mathrm{loss}}
% =
% \frac{N_{\mathrm{rej}}}{N_{\mathrm{arr}}},
% $$
% где $N_{\mathrm{arr}}$ --- число поступивших заявок, $N_{\mathrm{rej}}$ --- число отклонённых заявок.

% Пропускная способность определяется как число завершённых обслуживаний в единицу времени наблюдения:
% $$
% \widehat{\theta}
% =
% \frac{N_{\mathrm{completed}}}{T_{\mathrm{obs}}},
% $$
% где
% $$
% T_{\mathrm{obs}}=T_{\max}-T_{\mathrm{warm}}.
% $$
% Для достаточно длинного интервала наблюдения в loss-системе выполняется приближённая связь
% $$
% \widehat{\theta}\approx \lambda(1-\widehat{P}_{\mathrm{loss}}).
% $$

% Среднее число заявок в системе задаётся временным средним:
% $$
% \widehat{\mathbb{E}}[n]
% =
% \frac{1}{T_{\mathrm{obs}}}
% \int_{T_{\mathrm{warm}}}^{T_{\max}} n(t)\,dt.
% $$
% Средний занятый ресурс:
% $$
% \widehat{\mathbb{E}}[u]
% =
% \frac{1}{T_{\mathrm{obs}}}
% \int_{T_{\mathrm{warm}}}^{T_{\max}} u(t)\,dt.
% $$

% Оценка стационарного распределения числа заявок вычисляется через время пребывания системы в состоянии $k$:
% $$
% \widehat{\pi}(k)=\frac{T_k}{T_{\mathrm{obs}}},
% $$
% где $T_k$ --- суммарное время, в течение которого после разогрева в системе находилось $k$ заявок. Тогда
% $$
% \widehat{\mathbb{E}}[n]
% =
% \sum_{k=0}^{K} k\widehat{\pi}(k).
% $$

% Среднее время обслуживания оценивается по завершённым заявкам:
% $$
% \widehat{\mathbb{E}}[S]
% =
% \frac{1}{N_{\mathrm{completed}}}
% \sum_{i=1}^{N_{\mathrm{completed}}} S_i.
% $$
% Среднее время пребывания:
% $$
% \widehat{\mathbb{E}}[T_{\mathrm{soj}}]
% =
% \frac{1}{N_{\mathrm{completed}}}
% \sum_{i=1}^{N_{\mathrm{completed}}} T_{\mathrm{soj},i}.
% $$
% В рассматриваемой системе отсутствует очередь, поэтому для принятой заявки
% $$
% T_{\mathrm{soj},i}=S_i.
% $$
% Обе метрики сохраняются отдельно для совместимости с более общей постановкой.

% Чувствительность оценивается через отклонение метрик от базового сценария. Для вероятности отказа используется абсолютный сдвиг в процентных пунктах:
% $$
% \Delta P_{\mathrm{loss}}(s,s_0)
% =
% 100\cdot
% \left(
% \widehat{P}_{\mathrm{loss}}(s)
% -
% \widehat{P}_{\mathrm{loss}}(s_0)
% \right).
% $$
% Для положительных метрик используется относительный сдвиг:
% $$
% \delta_m(s,s_0)
% =
% 100\cdot
% \left(
% \frac{\widehat{m}(s)}{\widehat{m}(s_0)}
% -
% 1
% \right).
% $$

% При анализе входящего потока базовым считается пуассоновский arrival:
% $$
% a_0=\text{Poisson}.
% $$
% Условный эффект входящего потока при фиксированном workload имеет вид
% $$
% \Delta_A P_{\mathrm{loss}}(a\mid\omega,\lambda,\sigma)
% =
% 100\cdot
% \left(
% \widehat{P}_{\mathrm{loss}}(a,\omega,\lambda,\sigma)
% -
% \widehat{P}_{\mathrm{loss}}(a_0,\omega,\lambda,\sigma)
% \right).
% $$

% При анализе workload базовым считается детерминированное обслуживание:
% $$
% \omega_0=\text{Deterministic}.
% $$
% Условный эффект workload при фиксированном входящем потоке:
% $$
% \Delta_W P_{\mathrm{loss}}(\omega\mid a,\lambda,\sigma)
% =
% 100\cdot
% \left(
% \widehat{P}_{\mathrm{loss}}(a,\omega,\lambda,\sigma)
% -
% \widehat{P}_{\mathrm{loss}}(a,\omega_0,\lambda,\sigma)
% \right).
% $$

% Для совместного анализа используется базовая пара
% $$
% (a_0,\omega_0)=(\text{Poisson},\text{Deterministic}).
% $$
% Совместный сдвиг вероятности отказа:
% $$
% \Delta_{A,W}P_{\mathrm{loss}}(a,\omega,\lambda,\sigma)
% =
% 100\cdot
% \left(
% \widehat{P}_{\mathrm{loss}}(a,\omega,\lambda,\sigma)
% -
% \widehat{P}_{\mathrm{loss}}(a_0,\omega_0,\lambda,\sigma)
% \right).
% $$

% Для проверки неаддитивности влияния arrival и workload вводится взаимодействие факторов:
% $$
% I_m(a,\omega,\lambda,\sigma)
% =
% \widehat{m}(a,\omega,\lambda,\sigma)
% -
% \widehat{m}(a_0,\omega,\lambda,\sigma)
% -
% \widehat{m}(a,\omega_0,\lambda,\sigma)
% +
% \widehat{m}(a_0,\omega_0,\lambda,\sigma).
% $$
% Если $I_m\approx 0$, эффекты входящего потока и workload близки к аддитивным. Если $I_m$ заметно отличается от нуля, то влияние workload зависит от типа входящего потока.

% Статистическая надёжность оценок контролируется по независимым репликациям. Для фиксированного сценария $s$ и метрики $m$ получается выборка
% $$
% Y_{s,m}^{(1)},Y_{s,m}^{(2)},\dots,Y_{s,m}^{(R_{\mathrm{rep}})}.
% $$
% Точечная оценка:
% $$
% \overline{Y}_{s,m}
% =
% \frac{1}{R_{\mathrm{rep}}}
% \sum_{r=1}^{R_{\mathrm{rep}}}Y_{s,m}^{(r)}.
% $$
% Выборочная дисперсия:
% $$
% S_{s,m}^2
% =
% \frac{1}{R_{\mathrm{rep}}-1}
% \sum_{r=1}^{R_{\mathrm{rep}}}
% \left(
% Y_{s,m}^{(r)}-\overline{Y}_{s,m}
% \right)^2.
% $$
% Стандартная ошибка:
% $$
% \mathrm{SE}_{s,m}
% =
% \frac{S_{s,m}}{\sqrt{R_{\mathrm{rep}}}}.
% $$
% Доверительный интервал уровня $1-\alpha$:
% $$
% \overline{Y}_{s,m}
% \pm
% t_{1-\alpha/2,\,R_{\mathrm{rep}}-1}\cdot \mathrm{SE}_{s,m}.
% $$
% В работе используется уровень доверия $95\%$. Различие между сценариями считается содержательным, если оно существенно превышает внутрисценарный разброс по репликациям.

% \section{Алгоритм имитационного моделирования}
% \label{sec:simulation_algorithm}

% Имитационная модель строит траекторию $GI/G/N/N/R$ loss-системы как последовательность событий. Очередь отсутствует: поступившая заявка либо немедленно принимается, либо немедленно теряется. Это определяет всю динамику модели.

% На вход одной репликации подаётся сценарий
% $$
% s=(a,\omega,\lambda,\sigma,N,K,R,F_D)
% $$
% и seed генератора случайных чисел. Начальное состояние пустое:
% $$
% n(0)=0,\qquad u(0)=0.
% $$

% Генерируется последовательность межприходных интервалов $A_i$ с распределением, заданным типом входящего потока $a$. Моменты поступления:
% $$
% t_i=t_{i-1}+A_i.
% $$
% Для каждой заявки генерируются ресурсное требование $D_i\sim F_D$ и workload $W_i\sim F_W$. Время обслуживания:
% $$
% S_i=\frac{W_i}{\sigma}.
% $$

% В каждый момент моделирования выбирается ближайшее событие: поступление заявки или завершение обслуживания. Между соседними событиями состояние постоянно. Если текущее состояние на интервале равно $(n_j,u_j)$, а длина интервала равна $\Delta t_j$, то в накопители добавляются величины
% $$
% n_j\Delta t_j,\qquad u_j\Delta t_j.
% $$
% Если интервал пересекает момент $T_{\mathrm{warm}}$, учитывается только его часть после разогрева.

% При поступлении заявки проверяются два ограничения:
% $$
% n(t_i-)<K,\qquad u(t_i-)+D_i\leq R.
% $$
% Если оба ограничения выполнены, заявка принимается:
% $$
% n(t_i+)=n(t_i-)+1,
% \qquad
% u(t_i+)=u(t_i-)+D_i.
% $$
% Одновременно планируется событие завершения:
% $$
% c_i=t_i+S_i.
% $$
% Если хотя бы одно ограничение нарушено, заявка получает отказ. Состояние системы не меняется, но увеличиваются счётчики отказов. Причина отказа фиксируется отдельно: ограничение по вместимости, по серверу или по ресурсу.

% При завершении обслуживания заявка освобождает сервер и ресурс:
% $$
% n(c_i+)=n(c_i-)-1,
% \qquad
% u(c_i+)=u(c_i-)-D_i.
% $$

% После окончания репликации формируются счётчики:
% $$
% N_{\mathrm{arr}},\quad
% N_{\mathrm{acc}},\quad
% N_{\mathrm{rej}},\quad
% N_{\mathrm{completed}},
% $$
% накопленное ресурсное время, накопленное время по числу заявок и, при включённом сохранении состояний, вектор $T_k$. Эти величины преобразуются в метрики:
% $$
% \widehat{P}_{\mathrm{loss}}
% =
% \frac{N_{\mathrm{rej}}}{N_{\mathrm{arr}}},
% \qquad
% \widehat{\theta}
% =
% \frac{N_{\mathrm{completed}}}{T_{\mathrm{obs}}},
% $$
% $$
% \widehat{\mathbb{E}}[u]
% =
% \frac{\mathrm{resource\_time}}{T_{\mathrm{obs}}},
% \qquad
% \widehat{\pi}(k)
% =
% \frac{T_k}{T_{\mathrm{obs}}}.
% $$
% Среднее число заявок вычисляется либо как временное среднее по накопителю, либо через стационарные вероятности:
% $$
% \widehat{\mathbb{E}}[n]
% =
% \sum_{k=0}^{K} k\widehat{\pi}(k).
% $$

% Одна репликация даёт один набор оценок. Для сценария выполняются $R_{\mathrm{rep}}$ независимых репликаций с разными seed. Далее по выборке репликаций вычисляются средние значения, стандартные ошибки и доверительные интервалы, заданные в разделе~\ref{sec:metrics}. В программной реализации результат одной репликации соответствует структуре \texttt{RunSummary}, а агрегированный результат сценария --- статистике по всем \texttt{RunSummary} с одинаковым ключом сценария.

% Алгоритм имеет три уровня. Первый уровень --- событийная симуляция одной траектории. Второй уровень --- независимые репликации одного сценария. Третий уровень --- сравнение сценариев, различающихся типом входящего потока, распределением workload и уровнем нагрузки. Именно третий уровень используется в главе~\ref{ch:results_and_sensitivity} для анализа чувствительности.

\section{Структура программного комплекса и стадии эксперимента}
\label{sec:software_architecture}

Учитывая специфику имитируемой системы $GI/G/N/N/R$, каждая репликация может быть рассмотренна как независимый процесс, отсутствует буффер, т.е. нет необходимости хранить память об объектах в системе и отсутствует зависимость от состояния системы, было принято решение проводить вычислительный используя CUDA ядра видеокарты, каждое из которых может запускать 1 репликацию параллельно. Программный комплекс состоит из Python+Rust пайплайна, где первый валидирует и подает в систему значения параметров системы, отрисовывает графики в конце, а второй используется для производительности на тяжелых вычислениях.  

\begin{equation*}
\begin{aligned}
    \text{Подготовка: } & \texttt{values.py} \longrightarrow \texttt{experiment\_values.json} \\
    \text{Вычисления и анализ: } & \xrightarrow{\texttt{cargo run -- full}} \texttt{suite\_result.json} \\
    \text{Визуализация результатов: }     & \longrightarrow \texttt{plots.py}
\end{aligned}
\end{equation*}

Весь пайлайн запускается при помощи \texttt{launcher.py}, который организует порядок активации отдельных частей и передает нужную конфигурацию запуска cargo. Первый файл в цепочке - \texttt{values.py}, который задает исходную конфигурацию серии. 
\begin{table}[h!]
\centering
\caption{Базовые параметры конфигурации вычислительного эксперимента}
\label{tab:program_params}
\begin{tabularx}{\textwidth}{|l|l|X|}
\hline
\textbf{Группа} & \textbf{Параметр} & \textbf{Описание / Значение} \\ \hline
\multirow{3}{*}{Эксперимент} & \texttt{REPLICATIONS} & Число независимых повторов  \\ \cline{2-3} 
 & \texttt{MAX\_TIME} & Общее время моделирования в репликации \\ \cline{2-3} 
 & \texttt{WARMUP\_TIME} & Длина периода разогрева системы \\ \hline
\multirow{2}{*}{Система} & \texttt{SERVERS\_N} & Количество обслуживающих каналов (серверов) \\ \cline{2-3} 
 & \texttt{TOTAL\_R} & Суммарный доступный ресурс системы \\ \hline
\multirow{2}{*}{Нагрузка} & \texttt{ARR\_LEVELS} & Уровни интенсивности $\lambda$ \\ \cline{2-3} 
 & \texttt{SIGMA\_LEVELS} & Скорость обслуживания $\sigma$ \\ \hline
\multirow{2}{*}{Распределения} & \texttt{WORKLOAD} & Выбранные семейства распределений \\ \cline{2-3} 
 & \texttt{ARRIVAL} & Выбранные семейства распределений \\ \hline
\end{tabularx}
\end{table}

Во время запуска \texttt{values\_validation.py} проверяет значения на осмысленность и соответствие условиям эксперимента, после чего \texttt{export\_values.py} создает JSON-конфигурацию, который затем используется вычислительным модулем.

Rust пайплайн передает значения из файла \texttt{experiment\_values.json} в \texttt{ExperimentConfig}. Затем модуль \texttt{scenario\_grid.rs} строит сетку сценариев. Каждый сценарий является одной точкой в пространстве параметров:
$$
s = (a,\omega,\lambda,\sigma,N,R,F_D),
$$
где $a$ --- тип входящего потока, $\omega$ --- распределение workload, $\lambda$ --- средняя интенсивность поступления, $\sigma$ --- скорость обслуживания, $N$ --- число серверов, $R$ --- общий ресурс, $F_D$ --- распределение требования ресурса.

Пусть $A$ --- множество типов входящего потока, $\Omega$ --- множество распределений workload, $\Lambda$ --- набор уровней $\lambda$, $\Sigma$ --- набор уровней $\sigma$. Тогда в зависимости от выбранного режима строятся разные декартовы произведения параметров:
$$
\begin{aligned}
\text{base:} \quad
& \Lambda \times \Sigma, \\
\text{workload-sensitivity:} \quad
& \Omega \times \Lambda \times \Sigma, \\
\text{arrival-sensitivity:} \quad
& A \times \Lambda \times \Sigma, \\
\text{combined-sensitivity:} \quad
& \Omega \times A \times \Lambda \times \Sigma,
\end{aligned}
$$

Для всех сценариев формируется \texttt{ScenarioSpec}. В нём хранится имя сценария, технический ключ, численная характеристика входящего потока и workload, параметры системы и распределение ресурса. После построения сетки модуль \texttt{experiments.rs} разделяет сценарии в список независимых прогонов.

Один CUDA thread (поток) выполняет 1 репликацию одного сценария, например при 3 значениях $\lambda$, 4 видах распределения workload  и 4 виднах распределения входящего потока заявок мы получаем сетку из 48 потоков - $N_{\mathrm{sc}}$. Если задать число репликаций, например 20, то общее число потоков будет $48 \cdot 20 = 960$. Все они запускаются параллельно, каждый на отдельном CUDA ядре - $N_{\mathrm{runs}} = N_{\mathrm{sc}} \cdot R_{\mathrm{rep}}$. Это позволят получить большой прирост скорости запуска одного эксперимента в текущей постановке, так как даже непромышленные видеокарты позволят запустить эксперимент с группами параметров, который бы занял намного дольше на непромышленном процессоре. Однако, стоит учитывать, что даже один поток, превышающий число ядер видеокарты, удвоит время выполнения. Также сужается постановка задачи до $GI/G/N/N/R$, процессор для этого универсален.  

Вычислительный backend получает список \texttt{RunRequest}, группирует батчи и передает в CUDA-ядро. Массив параметров содержит: $\lambda, \sigma, T_{\max}, T_{\mathrm{warm}}, N, K, R,$ а также численные коды распределений arrival и workload, параметры гиперэкспоненциальных распределений, массив значений ресурса $D$ и соответствующая функция распределения. Внутри происходит имитация дискретно-событийной loss-системы, которая хранит текущее число заявок $n(t)$, занятый ресурс $u(t)$, параметры активных заявок и моменты их завершения. 

По завершении каждого прогона backend формирует \texttt{RunSummary}, в которой содержатся счётчики поступлений, принятых заявок, отказов, завершений, вероятность отказа $\widehat{P}_{\mathrm{loss}}$, $\widehat{\theta}$, $\widehat{\mathbb{E}}[u]$, $\widehat{\mathbb{E}}[n]$, среднее время обслуживания $\widehat{\mathbb{E}}[S]$ и среднее время пребывания $\widehat{\mathbb{E}}[T_{\mathrm{soj}}]$.

После завершения всех прогонов модуль \texttt{experiments.rs} группирует результаты по ключу сценария. Для каждого сценария получается выборка значений метрики по независимым репликациям:
$$
Y_{s,m}^{(1)},Y_{s,m}^{(2)},\ldots,Y_{s,m}^{(R_{\mathrm{rep}})}.
$$
Модуль \texttt{stats.rs} вычисляет выборочное среднее:
$$
\overline{Y}_{s,m}
=
\frac{1}{R_{\mathrm{rep}}}
\sum_{r=1}^{R_{\mathrm{rep}}}Y_{s,m}^{(r)},
$$
выборочное стандартное отклонение:
$$
S_{s,m}
=
\sqrt{
\frac{1}{R_{\mathrm{rep}}-1}
\sum_{r=1}^{R_{\mathrm{rep}}}
\left(
Y_{s,m}^{(r)}-\overline{Y}_{s,m}
\right)^2
},
$$
стандартную ошибку:
$$
\mathrm{SE}_{s,m}
=
\frac{S_{s,m}}{\sqrt{R_{\mathrm{rep}}}},
$$
и доверительный интервал:
$$
\overline{Y}_{s,m}
\pm
z_{1-\alpha/2}\mathrm{SE}_{s,m}.
$$

Модуль \texttt{output.rs} получает на вход значения \texttt{ExperimentSuiteResult}, а именно имя серии, время создания, уровень доверия и набор агрегированных по всем сценариям результатов, и сохраняет результаты в директорию вида $\texttt{results/<timestamp>\_\_<suite\_name>}$.
В неё записываются:
\begin{itemize}
    \item \texttt{suite\_result.json} --- полный результат серии, включая статистику по сценариям и отдельным репликациям;
    \item \texttt{aggregated\_summary.csv} --- плоская агрегированная таблица по сценариям;
    \item \texttt{metric\_summaries.csv} --- таблица статистик по метрикам;
    \item \texttt{run\_summaries.csv} --- результаты отдельных репликаций;
    \item \texttt{suite\_summary.txt} --- текстовая сводка серии.
\end{itemize}

После сохранения результатов \texttt{plots.py} читает \texttt{suite\_result.json}, преобразует вложенную структуру в таблицы сценариев и репликаций, после чего строит графики для анализа чувствительности. Файл \texttt{plots\_reworked.py} создан для анализа комбинированного сценария. 

Таким образом Python отвечает за параметры эксперимента и визуализацию, Rust --- за тяжелые вычисления, которые требуют строгой типизации параметров, построение сценариев, агрегацию и сохранение результатов, а GPU-backend --- за параллельное выполнение независимых репликаций по методу Монте-Карло. 

\section{Контроль корректности и статистической надёжности результатов}
\label{sec:verification}
% Текст раздела 2.6...